\chapter{代数：方程与不等式}\label{ch:g10:algebra}

代数是用字母进行计算的艺术。本章复习并深化
两种基本运算 --- 展开与因式分解 --- 并借助它们
系统地求解方程与不等式，以符号表
为核心工具。

\section{展开与因式分解}

\begin{definition}[展开、因式分解]\label{def:g10:algebra:expand}
\emph{展开}\index{展开}一个式子，是指把乘积化为和，
\emph{因式分解}\index{因式分解}则是其逆过程，
即把和化为乘积。基本法则是分配律：
\[
k(a + b) = ka + kb,
\qquad
(a + b)(c + d) = ac + ad + bc + bd .
\]
\end{definition}

\begin{example}\label{ex:g10:algebra:expand}
逐步展开 $(2x + 3)(x - 5)$：
\begin{align*}
(2x + 3)(x - 5)
&= 2x \times x + 2x \times (-5) + 3 \times x + 3 \times (-5) \\
&= 2x^2 - 10x + 3x - 15 \\
&= 2x^2 - 7x - 15 .
\end{align*}
\end{example}

\begin{theorem}[乘法公式]\label{thm:g10:algebra:identities}
对任意实数 $a$ 与 $b$，有：
\[
(a+b)^2 = a^2 + 2ab + b^2,
\qquad
(a-b)^2 = a^2 - 2ab + b^2,
\qquad
(a+b)(a-b) = a^2 - b^2 .
\]
\end{theorem}

\begin{proof}
每一式都可直接展开。对第一式：
\[
(a+b)^2 = (a+b)(a+b) = a^2 + ab + ba + b^2 = a^2 + 2ab + b^2 .
\]
第二式将 $b$ 换成 $-b$ 即得：$(a-b)^2 = a^2 + 2a(-b) +
(-b)^2 = a^2 - 2ab + b^2$。对第三式：
$(a+b)(a-b) = a^2 - ab + ba - b^2 = a^2 - b^2$。
\end{proof}

\begin{omfigure}
\begin{tikzpicture}[scale=0.62]
  \draw[thick] (0,0) rectangle (5,5);
  \draw[thick] (3.4,0) -- (3.4,5);
  \draw[thick] (0,3.4) -- (5,3.4);
  \fill[omDef!12] (0,0) rectangle (3.4,3.4);
  \fill[omThm!12] (3.4,3.4) rectangle (5,5);
  \fill[omProp!12] (3.4,0) rectangle (5,3.4);
  \fill[omProp!12] (0,3.4) rectangle (3.4,5);
  \node[font=\small] at (1.7,1.7) {$a^2$};
  \node[font=\small] at (4.2,4.2) {$b^2$};
  \node[font=\small] at (4.2,1.7) {$ab$};
  \node[font=\small] at (1.7,4.2) {$ab$};
  \node[below, font=\small] at (1.7,0) {$a$};
  \node[below, font=\small] at (4.2,0) {$b$};
  \node[left, font=\small] at (0,1.7) {$a$};
  \node[left, font=\small] at (0,4.2) {$b$};
\end{tikzpicture}

{\small 为何 $(a+b)^2 = a^2 + 2ab + b^2$：边长为 $a + b$ 的正方形
分成边长为 $a$ 的正方形、边长为 $b$ 的正方形，以及两个 $a \times b$
矩形。}
\end{omfigure}

\begin{example}\label{ex:g10:algebra:identities}
正用（展开）与逆用（因式分解）：
\begin{align*}
(3x + 2)^2 &= 9x^2 + 12x + 4,
& x^2 - 25 &= (x+5)(x-5), \\
(x - 4)^2 &= x^2 - 8x + 16,
& 4x^2 + 4x + 1 &= (2x + 1)^2 .
\end{align*}
这一方法也便于心算：$101 \times 99 = (100+1)(100-1) =
10000 - 1 = 9999$。
\end{example}

\begin{method}[式子的因式分解]\label{met:g10:algebra:factor}
按以下顺序尝试：
\begin{enumerate}
  \item \emph{提公因式：}找出每一项都含有的因式，并把它
        提到括号外，例如 $6x^2 + 4x = 2x(3x + 2)$。公因式也可以是
        整个括号，如 $(x+1)(2x-3) + (x+1)(x+7) = (x+1)(3x+4)$。
  \item \emph{乘法公式：}识别 $a^2 - b^2$，或
        $a^2 \pm 2ab + b^2$ 形式，例如 $9x^2 - 16 = (3x)^2 - 4^2 =
        (3x+4)(3x-4)$。
  \item 综合运用两种方法，并将结果展开以检验。
\end{enumerate}
\end{method}

\begin{example}\label{ex:g10:algebra:factorsteps}
因式分解 $E = (2x + 1)^2 - (x - 3)^2$。这是两个平方之差
$a^2 - b^2$ 的形式，其中 $a = 2x+1$，$b = x-3$：
\begin{align*}
E &= \bigl[(2x+1) + (x-3)\bigr] \times \bigl[(2x+1) - (x-3)\bigr] \\
  &= (3x - 2)(2x + 1 - x + 3) \\
  &= (3x - 2)(x + 4).
\end{align*}
注意第二个括号：减去 $x - 3$ 即减去 $x$
\emph{并}加上 $3$。
\end{example}

\section{方程}

\begin{definition}[方程、解]\label{def:g10:algebra:equation}
含未知数 $x$ 的等式称为\emph{方程}\index{方程}；
使等式成立的 $x$ 的值称为\emph{解}；
解方程即求出其\emph{全部}解。
\end{definition}

若两个方程的解完全相同，则称它们\emph{同解}。
两边同时加上同一个数，或两边同时乘以同一个
\emph{非零}数，所得方程与原方程同解。

\begin{example}[一元一次方程]\label{ex:g10:algebra:firstdegree}
解方程 $5x - 7 = 2x + 8$，每次只作一步变形：
\begin{align*}
5x - 7 &= 2x + 8 \\
5x - 2x &= 8 + 7 && \text{（两边同时加上 $7 - 2x$）} \\
3x &= 15 \\
x &= 5 && \text{（两边同时除以 $3$）。}
\end{align*}
方程的唯一解为 $5$。检验：$5 \times 5 - 7 = 18$，
$2 \times 5 + 8 = 18$。
\end{example}

\begin{theorem}[零乘积法则]\label{thm:g10:algebra:zeroproduct}
实数之积为零，当且仅当其中至少有一个
因式为零：
\[
A \times B = 0
\quad\Longleftrightarrow\quad
A = 0 \ \text{ 或 } \ B = 0 .
\]
\end{theorem}

\begin{proof}
若 $A = 0$ 或 $B = 0$，则显然乘积为 $0$。反之，设
$AB = 0$ 且 $A \neq 0$。两边除以 $A$，得 $B = 0$。
\end{proof}

\begin{method}[用因式分解法解方程]\label{met:g10:algebra:zeroproduct}
解右端为 $0$ 的方程时，
若其左端可以因式分解，则按以下步骤进行：
\begin{enumerate}
  \item 将所有项移到左边，使方程化为 $E = 0$；
  \item 将 $E$ 分解为若干一次因式的乘积；
  \item 应用零乘积法则，分别解各个小方程；
  \item 汇总所有求得的解。
\end{enumerate}
\end{method}

\begin{example}\label{ex:g10:algebra:zeroproduct}
解 $(x - 2)(3x + 6) = 0$：由 $x - 2 = 0$ 得 $x = 2$，
或由 $3x + 6 = 0$ 得 $x = -2$。方程的解为 $-2$ 和 $2$。

解 $x^2 = 9x$。注意\emph{不}要除以 $x$（它可能为零！），
而应把所有项移到一边并因式分解：
\[
x^2 - 9x = 0
\quad\Longleftrightarrow\quad
x(x - 9) = 0
\quad\Longleftrightarrow\quad
x = 0 \ \text{ 或 } \ x = 9 .
\]
\end{example}

\begin{example}[分式方程]\label{ex:g10:algebra:quotient}
解 $\dfrac{x - 1}{x + 2} = 0$。分式为零，当且仅当其
分子为零\emph{且}分母不为零。此时 $x - 1 = 0$
得 $x = 1$，且 $1 + 2 = 3 \neq 0$，故唯一解为 $1$。
$x = -2$ 是\emph{不允许}的值（使分母为零），
应始终首先排除。
\end{example}

\section{不等式与符号表}

\begin{proposition}[不等式的性质]\label{prop:g10:algebra:ineqrules}
设 $a \leq b$。
\begin{enumerate}
  \item 对任意 $c$，有 $a + c \leq b + c$（加法不改变不等号方向）。
  \item 若 $c > 0$，则 $ac \leq bc$（乘以正数，
        不等号方向不变）。
  \item 若 $c < 0$，则 $ac \geq bc$（乘以负数，
        \emph{改变}不等号方向）。
\end{enumerate}
\end{proposition}

\begin{proof}
1. $(b + c) - (a + c) = b - a \geq 0$。
2. $bc - ac = (b-a)c$ 为两个非负数之积，故
$bc \geq ac$。
3. 此时 $(b-a)c \leq 0$，因为它是非负数与负数之积，
故 $bc \leq ac$。
\end{proof}

\begin{example}\label{ex:g10:algebra:firstineq}
解不等式 $-2x + 3 > 9$：
\begin{align*}
-2x + 3 &> 9 \\
-2x &> 6 && \text{（两边减去 3）} \\
x &< -3 && \text{（两边除以 $-2 < 0$，改变不等号方向）。}
\end{align*}
解集为 $\intoo{-\infty}{-3}$。
\end{example}

\begin{definition}[符号表]\label{def:g10:algebra:signtable}
\emph{符号表}\index{符号表}按每个因式一行的方式，记录该表达式各因式的
符号（$+$ 或 $-$），即当 $x$ 取遍 $\R$ 时，
在每个使该因式为零的取值处记为 $0$；最后一行给出
乘积的符号，逐列应用符号法则确定。
\end{definition}

\begin{example}\label{ex:g10:algebra:signtable}
求 $P(x) = (x - 2)(3 - x)$ 的符号。因式 $x - 2$ 在 $2$ 处为零，之后为
正；因式 $3 - x$ 在 $3$ 处为零，之前为正：
\begin{center}
\renewcommand{\arraystretch}{1.3}
\begin{tabular}{c|ccccc}
$x$ & $-\infty$ \quad & $2$ & & $3$ & \quad $+\infty$ \\
\hline
$x - 2$ & $-$ & $0$ & $+$ & & $+$ \\
$3 - x$ & $+$ & & $+$ & $0$ & $-$ \\
\hline
$P(x)$ & $-$ & $0$ & $+$ & $0$ & $-$ \\
\end{tabular}
\end{center}
由此可知 $P(x) > 0$ 恰在 $\intoo{2}{3}$ 上成立，且 $P(x) \leq 0$ 在
$\intoc{-\infty}{2} \cup \intco{3}{+\infty}$ 上成立。
\end{example}

\begin{omfigure}
\begin{tikzpicture}
\begin{axis}[omaxis,
  width=8.2cm, height=5.2cm,
  xmin=-0.6, xmax=5.4, ymin=-2.6, ymax=1.4,
  xtick={2,3}, ytick=\empty]
  \addplot[omDef, thick, domain=0.35:4.65] {(x-2)*(3-x)};
  \fill (axis cs:2,0) circle (1.5pt);
  \fill (axis cs:3,0) circle (1.5pt);
  \node[omDef, font=\small] at (axis cs:2.5,0.65) {$P > 0$};
  \node[omThm, font=\small] at (axis cs:0.9,-1.3) {$P < 0$};
  \node[omThm, font=\small] at (axis cs:4.1,-1.3) {$P < 0$};
\end{axis}
\end{tikzpicture}

{\small 函数 $P(x) = (x-2)(3-x)$ 的图像印证了符号表：在根 $2$ 与 $3$ 之间
为正，在两根之外为负。}
\end{omfigure}

\begin{method}[用符号表解不等式]\label{met:g10:algebra:signtable}
求解形如 $E(x) \geq 0$ 或 $E(x) < 0$ 的不等式：
\begin{enumerate}
  \item 将所有项移至左边，使右边为
        $0$（切勿两边同乘符号未知的
        式子）；
  \item 将左边因式分解；若左边为分式，则
        通分。
  \item 按每个因式一行列出符号表（若为商，则在最后一行
        用双竖线标出禁止取值）；
  \item 读出所求符号成立的区间，并检验
        每个端点是否属于解集。
\end{enumerate}
\end{method}

\begin{example}\label{ex:g10:algebra:quotientineq}
解 $\dfrac{x + 1}{x - 3} \leq 0$。分子在 $-1$ 处为零，分母
在 $3$ 处为零（禁止值）。商在两者符号相反时为负，
这发生于 $x$ 介于 $-1$ 与 $3$ 之间时。
包含使\emph{分子}为零的端点，但排除
禁止值：解集为 $\intco{-1}{3}$。
\end{example}

\section{练习}

\begin{exercise}[$\star$]\label{exo:g10:algebra:1}
展开并化简：
\[
3(2x - 5) - 2(x - 7), \qquad
(x + 4)(2x - 1), \qquad
(3x - 2)^2, \qquad
(5 - x)(5 + x).
\]
\end{exercise}

\begin{exercise}[$\star$]\label{exo:g10:algebra:2}
因式分解：
\[
x^2 - 49, \qquad
x^2 + 10x + 25, \qquad
7x^2 - 14x, \qquad
16x^2 - 8x + 1 .
\]
\end{exercise}

\begin{exercise}[$\star$]\label{exo:g10:algebra:3}
解方程：
\[
4x + 9 = x - 3, \qquad
\frac{2x - 1}{3} = 5, \qquad
2(x - 3) = 2x + 1 .
\]
（其中有一个方程无解 --- 请说明理由。）
\end{exercise}

\begin{exercise}[$\star$]\label{exo:g10:algebra:4}
利用零乘积法则解方程：
\[
(x + 5)(2x - 8) = 0, \qquad
x^2 - 16 = 0, \qquad
x^2 + 6x + 9 = 0 .
\]
\end{exercise}

\begin{exercise}[$\star$]\label{exo:g10:algebra:5}
解下列不等式，并用区间表示解集：
\[
3x - 5 \leq 7, \qquad
-4x + 1 > 9, \qquad
2(x - 1) \geq 5x + 4 .
\]
\end{exercise}

\begin{exercise}[$\star\star$]\label{exo:g10:algebra:6}
       将 $E = (x - 1)(x + 3) - (x - 1)(2x - 5)$ 因式分解，然后解 $E = 0$。
\end{exercise}

\begin{exercise}[$\star\star$]\label{exo:g10:algebra:7}
       解方程 $x^2 = 5x$，再解 $(2x + 1)^2 = (x - 4)^2$。（对于第二个方程，
       将所有项移到左边，并对平方差进行因式分解。）
\end{exercise}

\begin{exercise}[$\star\star$]\label{exo:g10:algebra:8}
       利用符号表解下列不等式：
\[
(x + 2)(4 - x) > 0,
\qquad
\frac{2x - 6}{x + 1} \geq 0 .
\]
\end{exercise}

\begin{exercise}[$\star\star$]\label{exo:g10:algebra:9}
       某流媒体服务每月收费 $9$ 元，另加一次性注册费
       $15$ 元；竞争对手每月收费 $12$ 元，且不收注册费。经过多少个月，
       第一家服务的总费用才更低？列出并求解一个
       不等式。
\end{exercise}

\begin{exercise}[$\star\star$]\label{exo:g10:algebra:10}
       解方程 $\dfrac{x+3}{x-1} = 2$。（两边同乘 $x - 1$，并先排除禁止值；
       或者将所有项移到左边，在
       同一公分母下通分。）
\end{exercise}

\begin{exercise}[$\star\star\star$]\label{exo:g10:algebra:11}
       设 $n$ 为整数。证明 $(n + 1)^2 - n^2$ 总是奇数，并证明
       两个相邻奇数的平方之差总是
       $8$ 的倍数。
\end{exercise}

\section{周末问题：从巴比伦到黄金分割}

\begin{problem}[{周末问题 --- 和积问题，
以花拉子米的作图方式配方，以及
方程 $x^2 = x + 1$}]
\label{pb:g10:algebra:1}
       四千年前，巴比伦的书吏就在泥板上解过这道题：
       \emph{已知两数之和与之积，求这两个数}。他们的
       办法 --- 在两数的
       半和两侧对称地寻找两数 --- 至今仍能解开每个二次方程，其现代名称是
       \emph{配方}。本题将学习这一方法，
       看花拉子米如何真正地配出一个
       正方形，并把这一方法用于所有比例中最负盛名的
       那一个。

\medskip
\noindent\textbf{第一部分 --- 巴比伦问题。}
\begin{enumerate}
  \item 热身（\cref{thm:g10:algebra:identities}）：
        将 $x^2 - 10x + 25$ 因式分解，再将 $9x^2 - 49$ 因式分解，接着 ---
        结合恒等式与平方差 --- 将
        $x^2 + 6x + 9 - 4$ 因式分解。
  \item 分别解方程 $(x + 1)(x + 5) = 0$ 与
        $(3x - 7)(3x + 7) = 0$
        （\cref{thm:g10:algebra:zeroproduct}）。
  \item 书吏的问题：两数之和为 $10$，积为
        $21$。仿照巴比伦的做法，将两数写成 $5 - t$
        与 $5 + t$（关于半和对称），并求出 $t$，
        再求出这两个数。
  \item 方法相同：和为 $14$，积为 $33$。
  \item 一般方法：设和为 $s$、积为 $p$，证明
        这两个数为
        \[
        \frac s2 - t
        \quad\text{和}\quad
        \frac s2 + t,
        \qquad\text{其中}\quad
        t^2 = \left(\frac s2\right)^2 - p ,
        \]
        并写出 $s$ 与 $p$ 所需满足的条件，使该问题
        有解。
\end{enumerate}

\medskip
\noindent\textbf{第二部分 --- 配方。}
\begin{enumerate}[resume]
  \item 展开 $(x - a)(x - b)$，并由此断定：方程
        $x^2 - sx + p = 0$ 的解恰为
        和为 $s$、积为 $p$ 的那一对数。（书吏们其实
        一直在不自觉地解二次方程。）
  \item 验证 $x^2 - 10x + 21 = (x - 5)^2 - 4$，并
        求解方程 $x^2 - 10x + 21 = 0$，方法是将右边写成
        平方差并因式分解。与第 3 题作比较。
  \item 用配方法解方程 $x^2 + 6x - 7 = 0$。
  \item 对 $x^2 + 4x + 5$ 配方，并说明为什么
        方程 $x^2 + 4x + 5 = 0$ \emph{没有}实数
        解。归纳一般判别法：一旦将其写成
        $(x + h)^2 + k$ 的形式，二次式何时
        取值为零？
  \item 这个名称是字面意义上的。九世纪的巴格达，
        花拉子米把 $x^2 + 6x = 7$ 画成边长为 $x$ 的正方形，
        在它的两条边上各拼上一个 $3 \times x$ 的矩形，
        并用缺少的
        $3 \times 3$ 角完成\emph{配方}。画出图形，说明它如何
        把方程化为 $(x + 3)^2 = 16$ --- 然后求解。
\end{enumerate}

\medskip
\noindent\textbf{第三部分 --- 符号、表格与两个
古老的承诺。}
\begin{enumerate}[resume]
  \item 作出 $(x - 3)(x + 2)$ 的符号表
        （\cref{def:g10:algebra:signtable}），并解
        不等式 $(x - 3)(x + 2) < 0$。
  \item 解不等式 $x^2 \leq 10x - 21$（将所有项移到一边，
        并利用第 7 题的结果）。
  \item 利用符号表解不等式 $\dfrac{x - 1}{x + 2} \geq 0$，
        注意禁止值。
  \item 兑现一个古老的承诺：
        初中分册中狄多女王围栏问题里的农夫，用 $20$ 米的围栏，所围的
        $x \times (10 - x)$ 围地的面积满足
        \[
        A(x) = x(10 - x) = 25 - (x - 5)^2 .
        \]
        验证该恒等式，并由此读出一个 \emph{证明}，
        证明了六年级只能猜测的结论：面积的最大值，
        以及取得最大值的唯一 $x$。
  \item 另一题：证明 $x + \dfrac1x \geq 2$ 对
        每个 $x > 0$ 成立，且仅当 $x = 1$ 时取等号。（两边同乘
        $x$，整理后识别出一个平方 --- 这正是
        算术--几何平均不等式在
        初中分册中欧几里得关系周末问题的新装扮。）
\end{enumerate}

\medskip
\noindent\textbf{第 IV 部分 --- 黄金分割。}
将长为 $1$ 的线段截成较大的一段 $x$，
较小的一段为 $1 - x$，使得
\emph{全长比大段，犹如大段比小段}：
$\frac{1}{x} = \frac{x}{1 - x}$.
\begin{enumerate}[resume]
  \item 证明该条件可化为 $x^2 + x - 1 = 0$。
  \item 用配方法求解，并保留表示长度的那个根：
        $x = \frac{\sqrt5 - 1}{2} \approx 0.618$。
  \item \emph{黄金比}为 $\varphi = \frac1x = \frac{2}{\sqrt5 - 1}$。
        将分母有理化（分子分母同乘以 $\sqrt5 + 1$，
        此处用到一个恒等式），
        由此
        证明 $\varphi = \frac{\sqrt5 + 1}{2} \approx 1.618$。
  \item \emph{严格}验证 $\varphi$ 的
        两个神奇恒等式：
        \[
        \varphi^2 = \varphi + 1,
        \qquad
        \frac{1}{\varphi} = \varphi - 1 .
        \]
  \item 压轴：说明为什么 $\varphi$ 是无理数
        （初中分册中无理数的周末问题给出了
        关于 $\sqrt5$ 的关键事实），然后计算下列比值 $\frac53$、
        $\frac85$、$\frac{13}{8}$、$\frac{21}{13}$，它们都是
        连续的 赫马钱德拉–斐波那契 数之比
        （初中分册中的节奏计数周末问题），精确到三位小数。它们
        似乎在追逐什么？（证明它们确实能追上它，
        属于数列各章的内容）。
\end{enumerate}
\end{problem}
